\documentclass[10pt,a4paper]{amsart}
\usepackage[margin=19mm]{geometry}
\usepackage{amsmath,amssymb,mathtools}
\usepackage[scr=rsfs,cal=euler]{mathalfa}
\usepackage{xfrac}
\usepackage[unicode,hidelinks,bookmarks=false]{hyperref}
\newcommand{\ie}{i.e.\ }
\newcommand{\cf}{cf.\ }
\newcommand{\resp}{respectively\ }
\newcommand{\Subsection}{\subsection}
\providecommand{\nobreakdash}{\nobreak\hbox{-}}
\setlength{\emergencystretch}{2em}
\multlinegap0pt
\usepackage{amssymb}
\usepackage{booktabs}
\usepackage{mathtools}
\usepackage[shortlabels]{enumitem}
\newtheorem{theorem}{Theorem}
\newtheorem{lemma}{Lemma}
\usepackage{cleveref}
\crefname{theorem}{Theorem}{Theorems}
\crefname{lemma}{Lemma}{Lemmas}
\newcommand{\Aut}{\operatorname{Aut}}
\newcommand{\Sym}{\operatorname{Sym}}
\newcommand{\bK}[1]{\overleftrightarrow{K}_{#1}}
\newcommand{\can}{\operatorname{can}}
\newcommand{\dL}{\overrightarrow{\Lambda}_{3}}
\newcommand{\dP}{\overrightarrow{P}_{3}}
\newcommand{\dV}{\overrightarrow{V}_{3}}
\newcommand{\dunion}{\mathbin{\dot\cup}}
\title{Self-complementary completions of six-vertex digraphs}
\author{Xinan Dai, Wenhao Deng, Yingdong Shi, Tailin Wu, Yuchen Yang}
\date{Source: arXiv:2609.20231v1 (CC BY 4.0)}
\begin{document}
\maketitle

\noindent\emph{Source and licence.} Self-complementary completions of six-vertex digraphs. Version arXiv:2609.20231v1 (CC BY 4.0), distributed by arXiv under the Creative Commons Attribution 4.0 International licence, http://creativecommons.org/licenses/by/4.0/, as recorded in the arXiv licence metadata for this exact version and shown on the abstract page https://arxiv.org/abs/2609.20231v1. Full licence notice: \texttt{licenses/arxiv-2609.20231-CC-BY-4.0.txt}.\par\smallskip

\noindent\textit{Editorial scope.} Counted unit: the article's theorem thm:main-classification (source0.tex lines 111--125). The article's complete original proof of it is delivered with it (source0.tex lines 252--272, one \texttt{proof} environment of 21 lines, which the article places away from the statement). No mathematics is written, repaired or re-derived: the statement and the proof are the article's own lines, in the article's own notation, and no retained line is substituted or re-flowed.  Retained uncounted as the source-local context the proof needs: lines 145--155; lines 233--250.  Nothing was omitted from any retained span, so the package carries no omission record.  No retained line contains a citation, so the wrapper carries no bibliography and no bibliography entry.  No retained line contains a figure environment, an image-inclusion command or drawing code, so no image is delivered and the package declares no asset dependency.  Editorial formatting only, in addition: an AMS \texttt{amsart} preamble that restates the article's own declarations and macros used by the retained lines, namely the environments \texttt{theorem}, \texttt{lemma} on separate counters, together with the standard packages those lines need.  The retained environments print in this document as Theorem 1, Lemma 1 in order of appearance, whereas the article prints the same items with the numbering of its own full document.  The complete native source of the article as distributed in the arXiv e-print for this exact version is bundled unchanged as \texttt{sources/210850/source0.tex}.
\begin{lemma}[Alternating-orbit criterion]\label{lem:criterion}
A digraph \(D\) is completable if and only if there is a permutation \(\sigma\in\Sym(V(D))\) such that
\begin{enumerate}
\item every nontrivial cycle of \(\sigma\) has even length, and \(\sigma\) has at most one fixed point;
\item for every \(j\ge0\),
\begin{equation}\label{eq:odd-disjoint}
A(D)\cap \sigma^{2j+1}(A(D))=\varnothing.
\end{equation}
\end{enumerate}
Equivalently, on each \(\langle\sigma\rangle\)-orbit in \(\Omega(V(D))\), the arcs of \(D\) lie in at most one of the two alternating parity classes.
\end{lemma}

\begin{table}[ht]
\caption{The five isomorphism classes of eight-arc obstructions. Counts refer to a fixed labelled six-element vertex set.}
\label{tab:classes}
\centering
\small
\setlength{\tabcolsep}{4pt}
\begin{tabular}{@{}lrrrl@{}}
\toprule
Digraph & \(|\Aut(D)|\) & Labelled copies & Packing maps & Converse \\
\midrule
\(\bK{3}\dunion\dP\) & 6  & 120 & 36 & itself \\
\(\bK{3}\dunion\dV\) & 12 & 60  & 36 & \(\bK{3}\dunion\dL\) \\
\(\bK{3}\dunion\dL\) & 12 & 60  & 36 & \(\bK{3}\dunion\dV\) \\
\(\bK{2}\dunion T_4^+\) & 6 & 120 & 48 & \(\bK{2}\dunion T_4^-\) \\
\(\bK{2}\dunion T_4^-\) & 6 & 120 & 48 & \(\bK{2}\dunion T_4^+\) \\
\bottomrule
\end{tabular}
\end{table}

\begin{theorem}\label{thm:main-classification}
Up to isomorphism, the noncompletable six-vertex digraphs with eight arcs are
\[
\bK{3}\dunion\dP,
\qquad
\bK{3}\dunion\dV,
\qquad
\bK{3}\dunion\dL,
\qquad
\bK{2}\dunion T_4^+,
\qquad
\bK{2}\dunion T_4^-.
\]
There are \(480\) labelled examples. Every member of the list is arc-minimal, and the list reduces to three classes when converse digraphs are identified.
\end{theorem}

\begin{proof}[Proof of \cref{thm:main-classification}]
Fix \(V=\{0,1,2,3,4,5\}\) and order the \(30\) possible arcs lexicographically. For each admissible permutation \(\sigma\), decompose \(\Omega(V)\) into \(\langle\sigma\rangle\)-orbits and split each orbit into its two alternating classes. By \cref{lem:criterion}, an arc set \(M\subseteq\Omega(V)\) is completable if and only if, for at least one admissible \(\sigma\), it meets at most one parity class on every orbit.

This predicate was evaluated for all
\[
\binom{30}{8}=5{,}852{,}925
\]
eight-arc sets. Exactly \(480\) fail for every admissible permutation. For each rejected set \(M\), define
\[
\can(M)=\min_{\pi\in S_6}\pi(M)
\]
with respect to the fixed lexicographic bit order. Two labelled digraphs have the same canonical mask exactly when they are isomorphic. The \(480\) rejected masks yield five canonical masks, represented by the five digraphs in the theorem.

Their orbit sizes under \(S_6\) are
\[
120,\quad 60,\quad 60,\quad 120,\quad 120,
\]
which sum to \(480\) and agree independently with \(6!/|\Aut(D)|\) using the automorphism orders in \cref{tab:classes}. Taking converses fixes the \(\bK{3}\dunion\dP\) class and exchanges the two remaining pairs.

For each representative \(D\) and each arc \(e\in A(D)\), the computation also records an admissible permutation satisfying \cref{lem:criterion} for \(D-e\). Thus every one-arc deletion is completable. Since every spanning subdigraph of a completable digraph is again completable, each of the five obstructions is arc-minimal.
\end{proof}

\end{document}
